% Appendix C (writer: W2 of the restructure; plan: research/paper-review/restructure/PLAN.md). Proofs for sec:kinks; lem:kink:secant and prop:algstat:product are stated here.
% Sources: split/kink/notes-final.md Thms 2.1-2.2, 3.2, 3.4, Cor 3.3, Prop 3.5, Remark 3.2'; split/teacher/notes-final.md Lemma 1.1, Remark 1.2, Def 1.4, Lemma 1.5;
% split/verify-kink.md §5-§6; split/verify-teacher.md §2-§3 (g >= 0 under the hypotheses of (ii); c >= 0 in Lemma 1.5(c); the offset of Def 1.4); base app-algstat.tex (prop:algstat:product).
\section{Proofs for \texorpdfstring{\cref{sec:kinks}}{Section \ref*{sec:kinks}} (kinks)}\label{app:kinks}

\subsection{The hull is the graph of the convex minorant}\label{app:kinks:hull}

Let $\mathcal P$ be a point family. By \cref{lem:setup:hull}(v),(vi) and its proof, the lower boundary $\underline C$ of $C$ is $+\infty$ left of $\alpha_0$, linear between consecutive vertices and equal to $\rho_q$ right of a last vertex $v_q$; the vertices are points of the family, attained at exactly their lengths, and $\alpha_i\to\infty$ if there are infinitely many, since lengths are integers. $\underline C$ is convex and nonincreasing, since its slopes $-\sigma_1<-\sigma_2<\dots\le0$ increase. Every point $(\ell(\pi),L(\pi))$ lies in $C$, hence on or above $\underline C$, so $\underline C(\alpha)\le \underline C(\ell(\pi))\le L(\pi)$ whenever $\ell(\pi)\le\alpha$: $\underline C\le\rho$, hence $\underline C\le\breve\rho$; and $\rho(\alpha_i)=\rho_i$, since $\rho(\alpha_i)\ge \underline C(\alpha_i)$ and the program at $v_i$ has length $\alpha_i$. Conversely, a convex $f\le\rho$ has $f(\alpha_i)\le\rho_i$, lies below the chord $\underline C$ between consecutive vertices, and has $f\le\rho\le\rho_q=\underline C$ right of a last vertex. So $\breve\rho=\underline C$.

\subsection{The kink lemma, and lines against secants}\label{app:kinks:lemma}

\emph{Proof of \cref{lem:kink:main}.} Let $\pi_{\rm r}$ bound the point. (i) On $\mathcal N$, $J'_c\ge J_{\rm r}(c)+\eta_{\mathcal N}(c)$ by definition. For $\pi\in\mathcal M$, the right line and $\ell(\pi)\ge\ell_{\mathcal M}$ give
\[
J'_c(\pi)\ge c\,\ell(\pi)+L_{\rm r}-g-\tau(\ell(\pi)-\alpha_{\rm r})=(c-\tau)\ell(\pi)+\tau\alpha_{\rm r}+L_{\rm r}-g\ge J_{\rm r}(c)-\Delta(c),
\]
since $c>\tau$. If $\pi_{\rm r}\in\mathcal M$, the line at $\pi_{\rm r}$ gives $g\ge L_{\rm r}-L'(\pi_{\rm r})+\tau(\alpha_{\rm r}-\ell(\pi_{\rm r}))\ge0$, and then $\Delta(c)\ge g\ge0$ because $\ell_{\mathcal M}\le\alpha_{\rm r}$. Finally $\Phi(c)\ge\inf_\pi J'_c(\pi)$ because $L\ge L'$, and $\Phi(c)\le J_c(\pi_{\rm r})\le J_{\rm r}(c)$.

(ii) $J_c(\pi)\le\Phi(c)+\epsilon$ implies $J'_c(\pi)\le J_{\rm r}(c)+\epsilon$. A program with this property is not in $\mathcal N$, so it is in $\mathcal M$, and the line gives $(c-\tau)(\ell(\pi)-\alpha_{\rm r})\le g+\epsilon$. Since $J'_c(\pi_{\rm r})\le J_{\rm r}(c)$, also $\pi_{\rm r}\in\mathcal M$, so $g\ge0$ by (i). If $\ell(\pi)\ge\alpha_{\rm r}$, the line and the length bound give $L'(\pi)\ge L_{\rm r}-g-\tau(g+\epsilon)/(c-\tau)=L_{\rm r}-(cg+\tau\epsilon)/(c-\tau)$; if $\ell(\pi)<\alpha_{\rm r}$, the line gives $L'(\pi)\ge L_{\rm r}-g$, which is larger, since $(cg+\tau\epsilon)/(c-\tau)-g=\tau(g+\epsilon)/(c-\tau)\ge0$.

(iii) $\eta_{\mathcal N}$ is an infimum of the functions $c\mapsto c(\ell(\pi)-\alpha_{\rm r})+L'(\pi)-L_{\rm r}$, each affine or identically $+\infty$, and $\eta_{\mathcal N}>-\infty$ because $J'_c\ge0$. (iv) Rearrange $J'_c(\pi)\le J_{\rm r}(c)+\epsilon$ and use $\ell(\pi)\ge\ell_{\mathcal M}$. (v) For $\pi\in\mathcal N$ put $x:=\alpha_{\rm r}-\ell(\pi)$; then $J'_c(\pi)-J_{\rm r}(c)=L'(\pi)-L_{\rm r}-cx$. This is at least $\vartheta$ if $x\le0$; if $x>0$ it is at least $\max\{\sigma x,\vartheta\}-cx$, which is at least $(\sigma-c)x\ge(\sigma-c)\vartheta/\sigma$ when $\sigma x\ge\vartheta$ and more than $\vartheta-c\vartheta/\sigma$ otherwise; in all cases it is at least $\vartheta(1-c/\sigma)$. \qed

If every program that bounds the point lies in $\mathcal N$, then $\eta_{\mathcal N}(c)\le J'_c(\pi_{\rm r})-J_{\rm r}(c)\le0$, and (ii) is vacuous.

Every $\tau\ge0$ has a right line on every $\mathcal M$ with length floor $\ell_{\mathcal M}$, with $g=L_{\rm r}+\tau(\alpha_{\rm r}-\ell_{\mathcal M})$, because $L'\ge0$; so a slope alone says nothing, and what matters is the offset. The \emph{secant} of $\pi\in\mathcal M$ with $\ell(\pi)=\alpha_{\rm r}+x$, $x>0$, is $\varsigma(\pi):=(L_{\rm r}-L'(\pi))/x$. A line bounds how far to the right a near-optimal program can lie; the secant decides whether a program at the point itself survives.

\begin{lemma}[Lines against secants]\status{proved}\src{split/teacher Def 1.4, Lemma 1.5; split/verify-teacher \S3}\label{lem:kink:secant}
Let $(\alpha_{\rm r},L_{\rm r})$ be bounded and $\mathcal M\subseteq\mathcal P$.
(a) If $(\tau,g)$ is a right line on $\mathcal M$, every $\pi\in\mathcal M$ that is $x>0$ bits longer than $\alpha_{\rm r}$ has $\varsigma(\pi)\le\tau+g/x$.
(b) If $\pi_1\in\mathcal M$ is $x_1>0$ bits longer than $\alpha_{\rm r}$, every right line on $\mathcal M$ of slope $\tau$ has $g\ge(\varsigma(\pi_1)-\tau)x_1$.
(c) Let a program $\pi_0$ have $\ell(\pi_0)=\alpha_{\rm r}$ and $L'(\pi_0)=L_{\rm r}$. For $0\le c<s_+$ some program of $\mathcal M$ longer than $\alpha_{\rm r}$ has $J'_c<J'_c(\pi_0)$; for $c>s_+$ none has $J'_c\le J'_c(\pi_0)$.
\end{lemma}

\begin{proof}
(a) and (b) are the line at $\pi$ and at $\pi_1$, rearranged. (c) $J'_c(\pi)-J'_c(\pi_0)=cx+L'(\pi)-L_{\rm r}=x(c-\varsigma(\pi))$. If $0\le c<s_+$, then $s_+>0$, so some such $\pi$ has $\varsigma(\pi)>c$ and $J'_c(\pi)<J'_c(\pi_0)$; if $c>s_+$, every such $\pi$ has $\varsigma(\pi)<c$ and $J'_c(\pi)>J'_c(\pi_0)$.
\end{proof}

\subsection{Kinks on one data set}\label{app:kinks:data}

\emph{Proof of \cref{cor:kink:locate}.} (a) Let $\pi_-$ bound $(\alpha_-,L_-)$, so $J_\lambda(\hat\pi)\le\lambda\alpha_-+L_-$. If $\ell(\hat\pi)<\alpha_--w$, the kink gives $J_\lambda(\hat\pi)\ge\lambda\alpha_-+L_-+(\sigma-\lambda)(\alpha_--\ell(\hat\pi))>\lambda\alpha_-+L_-$ (for $\sigma=\infty$, $L(\hat\pi)=\infty$), a contradiction. The fit bound is $J_\lambda(\hat\pi)\le\lambda\alpha_-+L_-$ rearranged, and for $\lambda\ge1$ subtracting $(\lambda-1)\ell(\hat\pi)$ gives $\ell(\hat\pi)+L(\hat\pi)\le\alpha_-+L_-+(\lambda-1)(\alpha_--\ell(\hat\pi))\le\alpha_-+L_-+(\lambda-1)w$. (b) If $\ell(\hat\pi)>\alpha_++w_+$, then $J_\lambda(\hat\pi)\ge\lambda\alpha_++L_++(\lambda-\tau)(\ell(\hat\pi)-\alpha_+)>J_\lambda(\pi_+)$ for $\pi_+$ bounding $(\alpha_+,L_+)$.

(c) Programs shorter than $\alpha_{\rm r}$ have fit above $L_{\rm r}\ge\rho(\alpha_{\rm r})$, so $\rho(\alpha_{\rm r})$ is attained at length $\alpha_{\rm r}$. By (a) and (b), every $\lambda$-MAP with $\lambda\in(\tau,\sigma)$ has length $\alpha_{\rm r}$ and hence fit $\rho(\alpha_{\rm r})$, and $\lambda$-MAPs exist (\cref{lem:setup:basic}). The $\lambda$-MAP point is thus the same on a nonempty open interval, so it is a vertex $v_i$ (\cref{lem:setup:hull}(iii)). At $\lambda=\sigma_i$ ($i\ge1$) and at $\lambda=\sigma_{i+1}>0$ a neighbouring vertex is also a $\lambda$-MAP point (\cref{lem:setup:hull}(v)), so $(\tau,\sigma)\subseteq(\sigma_{i+1},\sigma_i)$. Conversely, by convexity $\breve\rho(\ell)\ge\rho_i+\sigma_i(\alpha_i-\ell)$ for $\ell<\alpha_i$ and $\breve\rho(\ell)\ge\rho_i-\sigma_{i+1}(\ell-\alpha_i)$ for $\ell>\alpha_i$, and every point lies on or above $\breve\rho$. \qed

\emph{The hypotheses are needed.} In the point family $\{(0,10),(1,5),(2,0)\}$ (one program at each point) the middle point is a zero-window $(5,5)$-kink; it is a $\lambda$-MAP point only at $\lambda=5$, together with both other points, and it is not a vertex, so (c) needs $\tau<\sigma$. In $\{(0,10),(2,0)\}$ the point $(2,0)$ is a zero-window left kink of steepness $5$, and at $\lambda=5$ the point $(0,10)$, of length $0<2$, is a $\lambda$-MAP; symmetrically, $(0,10)$ is a zero-window right kink of flatness $5$, and at $\lambda=5$ the point $(2,0)$ is a $\lambda$-MAP. So (a) needs $\lambda<\sigma$, and (b) needs $\lambda>\tau$.

\emph{Proof of \cref{cor:kink:split}.} (a) By \cref{lem:setup:coding} and $b$-sufficiency, $L(\pi;D)\ge K(y_D\mid\dom D)-\kappa-\ell(\pi)\ge L_{\rm r}-(\ell(\pi)-\alpha_{\rm r})-(b+\kappa)$ for every $\pi$: a right line $(1,b+\kappa)$ with length floor $0$. At $\pi_s$ it gives $b+\kappa\ge0$, hence the right kinks (the remark after \cref{def:kink:kinks}). For a $\lambda$-MAP $\hat\pi$, subtract $\ell(\hat\pi)+L(\hat\pi;D)\ge\alpha_{\rm r}+L_{\rm r}-(b+\kappa)$ from $\lambda\ell(\hat\pi)+L(\hat\pi;D)\le\lambda\alpha_{\rm r}+L_{\rm r}$ to get $(\lambda-1)(\ell(\hat\pi)-\alpha_{\rm r})\le b+\kappa$, and insert this into the line for the fit bound. (b) is (a) with \cref{cor:kink:locate}(a), whose two-part bound and $\alpha_{\rm r}+L_{\rm r}\le K(y_D\mid\dom D)+b$ give the sufficiency. \qed

At a program attaining the product sophistication, \cref{cor:kink:split}(a) gives the following.

\begin{proposition}[Steeper two-part codes over products]\status{proved}\src{steeper Prop 3.6; split/kink Cor 3.3}\label{prop:algstat:product}
Let $\lambda>1$. Every $\lambda$-MAP over $\Hcl{\times}$ for $D$ has length at most $\sophx_b(D)+(b+\kappa)/(\lambda-1)$ and is $\big(b+(\lambda-1)\sophx_b(D)\big)$-product-sufficient.
\end{proposition}

\begin{proof}
If $\sophx_b(D)=\infty$ there is nothing to prove; otherwise let $\pi_s$ attain it. \Cref{cor:kink:split}(a) in $\Hcl{\times}$ gives the length bound, and, since $\ell(\hat\pi)\ge0$, $\ell(\hat\pi)+L(\hat\pi;D)\le J_\lambda(\hat\pi;D)\le J_\lambda(\pi_s;D)\le K(y_D\mid\dom D)+b+(\lambda-1)\sophx_b(D)$.
\end{proof}

\emph{Proof of \cref{thm:kink:converse}.} The kinks give the lengths by \cref{cor:kink:locate}(a),(b) with zero windows. Conversely, assume the condition on lengths. The index $i$ exists, since $\sigma_0=\infty$ and the $\sigma_i$ decrease to $0$ or end with $\sigma_{i+1}=0$; for the same reason $j$ exists when $\tau>0$ or the vertices are finitely many. On the nonempty interval $(\max\{\tau,\sigma_{i+1}\},\sigma)\subseteq(\sigma_{i+1},\sigma_i)$ the $\lambda$-MAPs are the programs at $v_i$, so $\alpha_i\ge\alpha'$. Since $\sigma_{j+1}\le\tau<\sigma$, $j\ge i$, and on $(\tau,\min\{\sigma,\sigma_j\})\subseteq(\sigma_{j+1},\sigma_j)$ the $\lambda$-MAPs are the programs at $v_j$, so $\alpha_j\le\alpha''$. By \cref{cor:kink:locate}(c), $v_i$ is a zero-window left kink of steepness $\sigma_i\ge\sigma$, hence of steepness $\sigma$, and $v_j$ one of flatness $\sigma_{j+1}\le\tau$, hence of flatness $\tau$, at the points $(\alpha_i,\rho(\alpha_i))$ and $(\alpha_j,\rho(\alpha_j))$ (\cref{app:kinks:hull}). The MAP length is nonincreasing in $\lambda$ (\cref{lem:setup:hull}(ii)), which gives the minimum $\alpha_i$ and the maximum $\alpha_j$. If $\tau=0$ and there are infinitely many vertices, then $\sigma_i\to0$ and $\alpha_i\to\infty$: the $\lambda$-MAP lengths are unbounded as $\lambda\downarrow0$, and no point is a zero-window right kink of flatness $0$, since $\rho$ decreases beyond every length; so both conditions fail. \qed

\emph{Proof of \cref{prop:kink:universal}.} (i) Let $\pi_{\rm r}$ bound $(\alpha_{\rm r},L_{\rm r})$. The hypothesis at $\pi_{\rm r}$ gives $L(\pi_U)\le L_{\rm r}+\alpha_{\rm r}+\kappa_0<\infty$, and the kink gives $L(\pi_U)\ge L_{\rm r}+\sigma(\alpha_{\rm r}-\ell_U)$, since $\ell_U<\alpha_{\rm r}-w$. So $\sigma\le(\alpha_{\rm r}+\kappa_0)/(\alpha_{\rm r}-\ell_U)$, and $\sigma>1+\epsilon$ forces $\alpha_{\rm r}-\ell_U<(\ell_U+\kappa_0)/\epsilon$. (ii) A $\lambda$-MAP $\hat\pi$ has $\lambda\ell(\hat\pi)+L(\hat\pi)\le\lambda\ell_U+L(\pi_U)\le\lambda\ell_U+L(\hat\pi)+\ell(\hat\pi)+\kappa_0$, so $(\lambda-1)\ell(\hat\pi)\le\lambda\ell_U+\kappa_0$, as in \cref{thm:absorb:absorption}(a). The hypothesis at $\pi_U$ gives $\ell_U+\kappa_0\ge0$; so the bound is at most $\ell_U+(\ell_U+\kappa_0)/\epsilon$ when $\lambda-1\ge\epsilon$, and exceeds it when $\ell_U+\kappa_0>0$ and $\lambda-1<\epsilon$. \qed
